\section{Diagn\'osticos de la calibraci\'on secuencial de Black-Litterman}

La calibraci\'on secuencial se implementa en \texttt{calibracion\_secuencial\_bl.py}
y se valida mediante \texttt{validacion\_3h\_rolling\_bl.py}. A continuaci\'on
se detallan los diagn\'osticos de cada etapa.

\subsection{Stage 1: Selecci\'on de familia de \textit{view}}

Se evaluaron cuatro familias: desempleo macro, momentum general, momentum
top/bottom 1Y y momentum top market-cap 6M. Para cada una se ejecut\'o
Black-Litterman con configuraci\'on base y se midi\'o el Sharpe medio y delta
Sharpe en $\mathcal{H}_1$. La familia ganadora fue desempleo macro con Sharpe
medio 0.854 y delta Sharpe 0.159.

\subsection{Stage 2: Estructura de $P$}

Se probaron 9 combinaciones de tasa de desempleo asumida
$U \in \{3\%, 4\%, 5\%\}$ y sensibilidad sectorial
$\beta \in \{1.0, 1.5, 2.0\}$. La combinaci\'on ganadora fue $U = 4\%$ con
$\beta = 1.5$, reflejando un escenario de desempleo moderadamente bajo con
respuesta sectorial intermedia.

\subsection{Stage 3: Intensidad de $Q$}

Se escal\'o el vector $Q$ por factores $q_{\text{scale}} \in \{0.5, 1.0, 1.5,
2.0\}$. El factor 1.5 result\'o \'optimo, indicando que amplificar
moderadamente la se\~nal macroecon\'omica mejora el desempe\~no fuera de
muestra sin introducir inestabilidad excesiva.

\subsection{Stage 4: Confianza $\Omega$}

La confianza $c \in \{0.25, 0.35, 0.50, 0.75, 1.00\}$ controla la
incertidumbre asignada a la \textit{view}. Con $c = 0.50$, la matriz $\Omega$
refleja una incertidumbre intermedia: ni tan alta que diluya la \textit{view},
ni tan baja que ignore el prior de equilibrio.

\subsection{Stage 5: Par\'ametro $\tau$}

El par\'ametro $\tau \in \{0.05, 0.10, 0.20, 0.50, 1.00\}$ escala la
incertidumbre del prior. El valor \'optimo $\tau = 0.20$ es mayor que el
valor te\'orico sugerido por la literatura ($\tau \approx 1/T \approx 0.0004$
para 2\,500 observaciones), lo que es consistente con la pr\'actica de usar
$\tau$ como un hiperpar\'ametro de calibraci\'on m\'as que como una
estimaci\'on puntual de la incertidumbre del prior.

\subsection{Resultados de la validaci\'on rolling}

La validaci\'on rolling utiliz\'o 10 ventanas de calibraci\'on en el escenario
con pandemia y 3 en el escenario sin pandemia. Los resultados completos por
ventana y perfil se encuentran en
\texttt{05\_Validacion\_modelo/outputs/resumen\_ejecucion\_3h.md}.
